 Our results indicate that, overall, application of any of the evaluated denoising techniques substantially improves sensitivity without introducing substantial bias into signal estimates when used in conjunction with optimal-echo combination. 

The tensorial version of MPPCA (TMMPCA) \citep{olesenTensorDenoisingMultidimensional2023} was found the most effective strategy for improving signal quality, irrespective of the experimental paradigm, for both low resolution (2.4 and 2 mm isotropic voxels) and high resolution (1.5 mm isotropic voxels) datasets. TMPPCA relies on a higher-order singular value decomposition that determines the noise floor by jointly exploiting noise estimation within and across echoes for a given voxel patch \citep{olesenTensorDenoisingMultidimensional2023}. The performance of TMPPCA was followed by independently applying NORDIC \citep{vizioli2021lowering} to each echo dataset as a secondary option. In addition to the improvements in signal quality, TMPPCA offers a practical advantage with respect to NORDIC as it operates solely on magnitude data, whereas NORDIC requires both magnitude and phase information \citep{dowdle2023evaluating,vizioli2021lowering}.\\

\subsubsection{Main findings for resting-state data}

Across resting-state datasets acquired at multiple spatial and temporal resolutions, thermal denoising consistently improved the temporal signal-to-noise ratio (tSNR) of the optimally combined multi-echo fMRI signal \citep{posse1999enhancement, liuGeometricViewSignal2022a}. These findings are consistent with prior reports in single-echo fMRI \citep{fernandesMPPCADenoisingFMRI2023a, comby2023denoising, veraart2016denoising}, and complements recent evidence in multi-echo fMRI data for precision functional mapping \citep{moserMultiechoAcquisitionThermal2025, krimmelHumanBrainstemsRed2025}. A consistent ordering of methods emerged across resolutions, although its structure depended systematically on voxel size. For  2.4 mm and 2 mm isotropic resolution, tSNR followed the hierarchy TMPPCA $>$ NORDIC  $\approx$ HYDRA $>$ MPPCA $>$ no denoising (VANILLA), whereas at higher spatial resolution (1.5 mm isotropic) the ordering changed slightly to TMPPCA $>$ NORDIC $>$ HYDRA $\approx$ MPPCA $>$ no denoising (VANILLA). These results can be interpreted in terms of how each method exploits redundancy in multidimensional (multicontrast) fMRI data. TMPPCA, which extends MPPCA to a tensor-based formulation \citep{olesenTensorDenoisingMultidimensional2023}, consistently provided the most effective separation of signal and thermal noise by jointly modeling structure across spatial, temporal and echo dimensions. NORDIC also improves denoising performance through explicit modeling of thermal noise in the complex domain, but does not fully exploit higher-order dependencies across echoes. In contrast, HYDRA, which was developed here as the application of NORDIC after spatial concatenation of multi-echo data, partially disrupts the intrinsic multi-echo structure, reducing its ability to leverage echo-dependent redundancy, which explains its convergence toward MPPCA performance at higher spatial resolution.

Furthermore, we evaluated whether multi-echo thermal denoising alters estimates of the net magnetization $\overline{S_0}$, and the apparent transverse relaxation rate $\overline{R_2^*}$ (or $\overline{T_2^*}$), which are the key parameters underlying the optimal combination method of echoes through $T_2^*$-weighted signal modelling \citep{posse1999enhancement, kundu2017multi, liuGeometricViewSignal2022a}. To our knowledge, this constitutes the first systematic evaluation of how thermal denoising affect voxel-wise parameter estimates in multi-echo fMRI. 

Overall, no substantial differences were observed in either $\overline{S_0}$ or $\overline{R_2^*}$ across thermal denoising methods at the whole brain level. However, small but systematic deviations emerged in regions characterised by low SNR and strong susceptibility effects, including such as orbitofrontal cortex, anterior frontal cortex, inferior temporal cortex, and subcortical areas. These regions are known to be particularly sensitive to field inhomogeneities, signal dropout, g-ratio noise amplication, and thus SNR degradation, induced by accelerated acquisitions (GRAPPA parallel imaging, partial Fourier, and simultaneous multislice acquisition). 

Within these brain regions, MPPCA and TMPPCA consistently led to a slight underestimation of $\overline{S_0}$ in the three datasets. Furthermore, for the lower resolution datasets (2.4 and 2 mm), MPPCA and TMPPCA also exhibited a modest underestimation of $\overline{R_2^*}$ in these regions. Jointly with the $S_0$ behaviour, this pattern suggests a reduction in estimated BOLD signal variance across echoes, effectively resulting in a shallower apparent echo signal decay with echo time. 
In contrast, NORDIC showed a mild overestimation of $\overline{S_0}$ and $\overline{R_2^*}$ in these low-SNR areas, indicating that some degree of signal “leakage” may occur during denoising, and thus some BOLD signal is removed. This behaviour is consistent with previous observations in single-echo fMRI denoising studies, where strong thermal noise suppression can induce regionally specific biases due to signal leakage in low-SNR areas, particularly when low-rank assumptions are applied globally across the brain \citep{fernandesMPPCADenoisingFMRI2023a, zhu2022denoise, faesEvaluatingEffectDenoising2024}. Such effects may reflect partial redistribution of variance across voxels induced by global low-rank approximation, potentially propagating signal from high SNR to low SNR voxels in a consistent manner across individual echoes and thus impacting the signal model. 

Further analyses of instantaneous changes in $R_2^*$, i.e. $\Delta R_2^*$, demonstrated consistent effects across all datasets, showing that multi-echo thermal denoising enhances physiologically plausible BOLD-related fluctuations \citep{caballerogaudesParadigmFreeMapping2011}. Similar to the tSNR analysis, the ordering of methods followed a stable hierarchy of more physiologically-plausible estimates: TMPPCA $>$ NORDIC  $\approx$ HYDRA $>$ MPPCA $>$ no denoising. These results indicate that the residual signal after multi-echo thermal denoising exhibits enhanced sensitivity to BOLD-related fluctuations. Importantly, this suggests that although thermal denoising has minimal impact on mean relaxation parameters at the global level, it improves improves the detectability of functionally relevant temporal fluctuations in multi-echo fMRI data. The fact that TMMPCA showed the superior performance in this context might improve previous studies of precision functional mapping with multi-echo fMRI, which either applied NORDIC independently to each echo \cite{moserMultiechoAcquisitionThermal2025, krimmelHumanBrainstemsRed2025} or no thermal denoising was applied \citep{lynchRapidPrecisionFunctional2020, lynchFrontostriatalSalienceNetwork2024}.

Reliability analysis of functional connectivity patterns \citep{lynchRapidPrecisionFunctional2020} demonstrated increased test–retest stability following multi-echo thermal denoising at lower spatial resolutions (2.4 mm and 2 mm isotropic voxels). In contrast, at higher spatial resolution, some voxels exhibited improved reliability while others showed decreased stability. Overall, the comparative performance of the denoising approaches can be summarized hierarchically as follows: TMPPCA $>$ NORDIC $\approx$ HYDRA $>$ MPPCA $>$ no denoising (VANILLA).
The use of multi-echo thermal denoising may enable shorter scan durations for FC fingerprinting or network identification \citep{finnFunctionalConnectomeFingerprinting2015, foxSpontaneousFluctuationsBrain2007}, boosting the advantages of using multi-echo fMRI for these applications \citep{lynchRapidPrecisionFunctional2020} thus as.

\subsubsection{Main findings for task-related data}

Given the central role of fMRI in mapping task-evoked brain responses in cognitive neuroscience, we conducted a task-based analysis in Chap. \ref{chap:taskdenoising} \citep{bermanEvaluatingFunctionalLocalizers2010, saxeDivideConquerDefense2006}. Previous studies of thermal denoising have also focused on the analysis of task-related activity, but all of them used single-echo fMRI data \citep{vizioli2021lowering, dowdle2023evaluating, }. These analysis have not been made before exploring the interaction between multi-echo and thermal denoising for task-based activations.  We employed a visual localizer specifically designed to identify activation in the fusiform face area (FFA) \citep{lerma-usabiagaConvergingEvidenceFunctional2018}. Under a controlled experimental paradigm, the observed brain activity is expected to be strongly constrained by the design matrix. Building on this assumption, we implemented a cross-validation framework to compare the fit of test time series as predicted by their corresponding design matrices \citep{princeImprovingAccuracySingletrial2022, dowdle2023evaluating}. 

Overall, model goodness-of-fit increased as a consequence of thermal denoising. Consistent with the resting-state analysis, we observed that improvements in goodness-of-fit propagated from regions with higher signal-to-noise ratio (SNR) located closer to the head coils toward deeper and more distal brain regions. This spatial pattern is expected, as g-factor noise amplification is higher in these areas \citep{aja-fernandezNoiseEstimationParallel2014, pruessmann1999sense}, thereby underscoring the efficacy of thermal denoising approaches in enhancing SNR particularly where g-factor and, consequently, thermal noise are more pronounced. However, the magnitude of the improvement in goodness-of-fit systematically followed the ordering: vanilla $<$ MPPCA $<$ HYDRA $<$ NORDIC $<$ TMPPCA. 

Subsequently, we performed a more conventional analysis by evaluating our ability to detect the FFA \citep{kanwisherFusiformFaceArea1997} across denoising strategies, using a minimal number of runs to fit the GLM (i.e., one, two, or three runs). We found that thermal denoising reduced the number of runs required to obtain a robust GLM fit while substantially improving task detection sensitivity. In this context, TMPPCA exhibited the best performance, achieving reliable FFA detection with a single run. NORDIC also enhanced sensitivity relative to the non-denoised data, although to a lesser extent than TMPPCA.\\


\section{Limitations and future work}

A key advantage of multi-echo acquisitions—beyond optimal combination—is the removal of non‑BOLD signal components using ME‑ICA \citep{kundu2012differentiating,kundu2017multi}. In the present study, we did not investigate the effect of thermal denoising on ICA decomposition or on subsequent ME‑ICA denoising, which remains an avenue for future research enabled by software such as tedana \citep{dupre2021te}. Future studies could examine how prior denoising affects the number of components estimated, the proportion classified as BOLD versus noise, and the overall decomposition quality. Future efforts could explore modifications to the ME‑ICA denoising pipeline that would improve the separation of BOLD and non‑BOLD components following thermal denoising. For example, the tedana workflow performs an initial PCA step specifically designed to reduce thermal noise \citep{dupre2021te}. If thermal noise is already suppressed before entering the tedana pipeline, this step may need to be modified or omitted, or at the very least, its effect on already‑denoised data should be systematically evaluated. 

Another area of interest is extending TMPPCA (or more generally, tensorial denoising using higher‑order singular value decomposition, HOSVD) to operate directly on complex‑valued fMRI data, rather than on magnitude images. In the complex domain, thermal noise follows a Gaussian distribution, whereas magnitude data converts it into a Rician distribution \citep{gudbjartssonRicianDistributionNoisy1995c}. This Rician floor biases low‑signal regions and introduces signal‑dependent noise, which violates the Gaussian assumptions underlying many denoising algorithms \citep{veraart2016denoising, comby2023denoising}. By performing HOSVD‑based tensor denoising on complex data, the noise remains additive, zero‑mean, and Gaussian, allowing more effective separation of signal and noise components. This, in turn, should improve the preservation of small BOLD fluctuations, reduce bias in parameter estimates, and increase overall denoising power \citep{wangModifiedHigherorderSingular2020}.



